Consider a stock that trades at $S = \$100$ today. A \emph{call option} on this stock
gives its holder the right, but not the obligation, to buy the stock at a fixed
\emph{strike} price, say $K = \$100$, at a fixed \emph{maturity}, say $T = 1$ year from
now. If the stock trades at \$130 in a year, the holder buys at \$100 and gains \$30;
if it trades at \$80, the holder simply does not use the option and loses nothing more.

This right has a price $C$ today, which is much smaller than the stock price: with
$\sigma = 0.2$ and $r = 0.02$ (defined below), $C \approx \$8.9$, i.e.\ $C/S \approx 0.089$.
The price depends strongly on how much the stock is expected to fluctuate, its
\emph{volatility} $\sigma$: if the stock is twice as volatile ($\sigma = 0.4$), large gains become
more likely while losses remain capped, and the same option costs $C \approx \$16.7$.